\documentclass{article}
\usepackage[a4paper]{geometry}
\usepackage{iftex}
\ifPDFTeX
  \usepackage[T1]{fontenc}
  \usepackage[utf8]{inputenc}
\else
  \usepackage{fontspec}
\fi
\usepackage[normalem]{ulem}% for \sout in the alternatives tests
\usepackage[lazy,gb4e]{linguexx}
\usepackage{hyperref}

\begin{document}

\section*{Numbering and nesting}

\ex. Plain example one.

\ex.
\a. First sub.
\b. Second sub.
\a. Sub-sub roman one.
\b. Sub-sub roman two.

\ex.\label{lbl:main} Labelled example: text continues.

\ex.[(4$'$)] Custom label, counter not stepped.

\ex. Counter continues correctly at four.

\section*{Judgments}

\ex. *Auto-detected star.

\ex. ??This one is doubly questionable.

\ex. ?*Mixed marks.

\ex. \#Hash via control sequence.

\ex. \%\#Percent-hash combination.

\ex.
\a. *Starred sub-example,
\b. plain sub-example: texts must align.

\ex. \jdg{\dag}Explicit arbitrary mark via jdg.

\section*{Glossing}

\ex. \gll Dies ist ein kleines Beispiel \\
     this is a small example \\
\glt `This is a small example.'

\ex. \gll Das {kleine Kind} schl\"aft \\
     the {little child} sleeps \\
\glt `The little child sleeps.' (braced group = one column)

\ex. \gll Eins zwei drei vier \\
     one two \\
\glt Unequal lines: empty gloss cells.

\ex. \glll Ich habe geschlafen \\
      ich hab-e schlaf-en \\
      1SG have-1SG sleep-PTCP \\
\glt `I slept.' (three-line gloss)

\ex. \gll Ein sehr sehr sehr langes Beispiel damit die Glossenzeilen
umgebrochen werden m\"ussen und wir sehen ob der Umbruch zwischen den
Spalten funktioniert wie er soll \\
one very very very long example so.that the gloss.lines broken
become must and we see whether the break between the columns
works as it should \\
\glt `Wrapping test.'

\GlossTierFont{4}{\textit}
\ex. \gl Der Hund schl\"aft \\
     der Hund schlaf-t \\
     DEF.NOM.SG.M dog sleep-3SG \\
     D N V \\ \endgl
\glt `The dog sleeps.' (four tiers, tier 4 italic via GlossTierFont)

\ex. \gl Unequal tiers and {braced pairs} here \\
     only two \\
     third {tier has} a braced group \\ \endgl
\glt Multi-tier: empty cells and braces.

\section*{References}

\ex.\label{r:one}
\a.\sublabel{r:a} Premier.
\b. Deuxi\`eme.
\c.\sublabel{r:c} Troisi\`eme.

Refs: \ref{r:one}, \ref{r:a}, \pref{r:c}; range \refrange{r:a}{r:c};
bare \prefrange{r:a}{r:c}; relative: \Last, \pLast, \LLast; and the next
example will be \Next.

\ex. As announced.

\section*{Footnote examples}

Main text example numbering\footnote{A footnote with examples:
\ex. First footnote example, roman.

\ex. *Second, with judgment.

Footnote-internal refs: \Last\ and \LLast.} is not disturbed:

\ex. Text example after the footnote; this is \Last.

\section*{Alternatives and sources}

\ex. \altn{\sout{This girl in the red coat}}{she}{Mary} will put a
picture of \altn{Bill}{him} on your desk. \exsource{(Perlmutter 1971)}

\ex.
\a. Multipart source A. \exsource{(Source A)}
\b. Multipart source B, same number. \exsource{(Source B)}

\ex. \gll Dies ist ein Beispiel mit Glosse und Quelle \\
     this is an example with gloss and source \\
\glt `Glossed with source.' \exsource{(Schaden 2020: 45)}

A very long example sentence follows to test the flush-right fallback.

\ex. This example sentence is specifically designed to run almost to the
end of the line so that the source cannot fit and must be pushed onto its
own line, flush right. \exsource{(Tolsto\"i 1869: t.~2, p.~334)}

\section*{Glossed shorthands and early termination}

\exg. *Das Beispiel funktionieren \\
      the example work.INF \\
\glt Judged glossed shorthand: star hangs left of the first column.

\ex.
\ag. Erstes Unterbeispiel mit Glosse \\
     first sub.example with gloss \\
\bg. Zweites auch \\
     second too \\

\ex. An example ended early by z. \z. This continuation must sit flush
left at the outer margin, outside the example.

\ex. Numbering continues undisturbed after z.

\ex.
\a. Sub level,
\a. sub-sub level, then \z.\z. two pops escape from two levels deep.

\ex. And the counter is still right.

\section*{Termination without blank lines (v3.2)}

\ex. Ended by z with no blank line after it. \z.This prose follows
immediately in the source, on the same line.
\ex. And this example follows with no blank line before it either.

\begin{minipage}{0.8\textwidth}
\ex. In a minipage, terminated by the environment boundary itself.\end{minipage}

\ex. With \begin{tabular}{ll} a & b \\ c & d \end{tabular} inside,
ended by a blank line: environment counting works.

{\ex. Inside a brace group, ended by the group's closing brace.}

Footnote examples without blank lines:\footnote{Before.
\ex. Footnote example ended by z. \z.After, still in the footnote.}
main text continues.

\ex.
\a. First
\b. Second
\a. Second a
\b. Second b
\z.
\b. third: z popped one level, so this is item c.
\z.
and this is prose outside the example, flush left.

\ex.
\a. A forgotten blank line after a sub-item:
\ex. this example is still set at top level, not misindented.

\section*{Mixing environments and dot syntax}

\ex.
\a. Dot sub-example,
\begin{xlist}
\ex with an xlist sub-sub inside,
\item made by ex or by plain item,
\end{xlist}
\b. and letters resume after it.

\section*{gb4e-compatible interface}

\begin{exe}
\ex\label{gb:one} First in a batch.
\ex[*]{Judged via the bracket form.}
\ex Plain, with nesting:
 \begin{xlist}
 \ex[]{Empty judgment,}
 \ex[??]{double mark,}
 \ex
  \begin{xlist}
  \ex[**]{roman with stars,}
  \ex[\%]{and a percent.}
  \end{xlist}
 \ex after the nesting.
 \end{xlist}
\ex Last of the batch; ref: \ref{gb:one}.
\end{exe}

\ex. The very last example, and the counter survived it all: this is \Last.

Final paragraph: normal text after everything, correctly outdented.

\end{document}
